\chapter{Septic Abortion}
\index{Septic Abortion}

\medskip
\noindent\textit{\textbf{Under review.} This chapter is an AI-assisted educational draft; it is not a substitute for professional medical advice.}

\section*{Definition and Overview}
Septic abortion is infection of the uterus and surrounding tissues after miscarriage or abortion, often with retained tissue. Fever, chills, pelvic pain, foul discharge, heavy bleeding, dizziness, or collapse require urgent emergency assessment. Treatment includes antibiotics, fluids, sepsis care, and uterine evacuation when indicated.

\section*{Clinical Features}
Symptoms and concerns vary with the specific topic, underlying cause, age, pregnancy status, and other health conditions.

\section*{When to Seek Medical Care}
Seek medical review for new, persistent, worsening, or function-limiting symptoms. Contact your local emergency services for severe breathing difficulty, collapse, sudden neurological change, severe chest or abdominal pain, or rapidly worsening illness.

\section*{Diagnosis and Investigations}
Assessment starts with history and examination, followed by targeted laboratory tests, imaging, monitoring, or specialist referral when indicated.

\section*{Management}
Management is individualised and may include education, self-care, medicines, procedures, rehabilitation, monitoring, and referral. Do not start, stop, or share prescription treatment without clinical advice.

\section*{Complications}
Complications depend on the underlying cause and may include progression, effects on other organs, treatment adverse effects, or reduced quality of life.

\section*{Prognosis and Self-Care}
Follow the care plan, keep appointments, avoid known triggers, and seek help if symptoms worsen. Individual prognosis should be discussed with a clinician.
\section*{Safety-netting}
Seek urgent help for severe or rapidly worsening symptoms, collapse, breathing difficulty, severe pain, confusion, dehydration, uncontrolled bleeding, or any immediate danger.
